% Options for packages loaded elsewhere
\PassOptionsToPackage{unicode}{hyperref}
\PassOptionsToPackage{hyphens}{url}
%
\documentclass[
]{book}
\usepackage{amsmath,amssymb}
\usepackage{iftex}
\ifPDFTeX
  \usepackage[T1]{fontenc}
  \usepackage[utf8]{inputenc}
  \usepackage{textcomp} % provide euro and other symbols
\else % if luatex or xetex
  \usepackage{unicode-math} % this also loads fontspec
  \defaultfontfeatures{Scale=MatchLowercase}
  \defaultfontfeatures[\rmfamily]{Ligatures=TeX,Scale=1}
\fi
\usepackage{lmodern}
\ifPDFTeX\else
  % xetex/luatex font selection
\fi
% Use upquote if available, for straight quotes in verbatim environments
\IfFileExists{upquote.sty}{\usepackage{upquote}}{}
\IfFileExists{microtype.sty}{% use microtype if available
  \usepackage[]{microtype}
  \UseMicrotypeSet[protrusion]{basicmath} % disable protrusion for tt fonts
}{}
\makeatletter
\@ifundefined{KOMAClassName}{% if non-KOMA class
  \IfFileExists{parskip.sty}{%
    \usepackage{parskip}
  }{% else
    \setlength{\parindent}{0pt}
    \setlength{\parskip}{6pt plus 2pt minus 1pt}}
}{% if KOMA class
  \KOMAoptions{parskip=half}}
\makeatother
\usepackage{xcolor}
\usepackage{longtable,booktabs,array}
\usepackage{calc} % for calculating minipage widths
% Correct order of tables after \paragraph or \subparagraph
\usepackage{etoolbox}
\makeatletter
\patchcmd\longtable{\par}{\if@noskipsec\mbox{}\fi\par}{}{}
\makeatother
% Allow footnotes in longtable head/foot
\IfFileExists{footnotehyper.sty}{\usepackage{footnotehyper}}{\usepackage{footnote}}
\makesavenoteenv{longtable}
\usepackage{graphicx}
\makeatletter
\def\maxwidth{\ifdim\Gin@nat@width>\linewidth\linewidth\else\Gin@nat@width\fi}
\def\maxheight{\ifdim\Gin@nat@height>\textheight\textheight\else\Gin@nat@height\fi}
\makeatother
% Scale images if necessary, so that they will not overflow the page
% margins by default, and it is still possible to overwrite the defaults
% using explicit options in \includegraphics[width, height, ...]{}
\setkeys{Gin}{width=\maxwidth,height=\maxheight,keepaspectratio}
% Set default figure placement to htbp
\makeatletter
\def\fps@figure{htbp}
\makeatother
\setlength{\emergencystretch}{3em} % prevent overfull lines
\providecommand{\tightlist}{%
  \setlength{\itemsep}{0pt}\setlength{\parskip}{0pt}}
\setcounter{secnumdepth}{5}
\usepackage{booktabs}
\usepackage{amsthm}
\makeatletter
\def\thm@space@setup{%
  \thm@preskip=8pt plus 2pt minus 4pt
  \thm@postskip=\thm@preskip
}
\makeatother

\usepackage{tcolorbox}


\newtcolorbox{blackbox}{
  colback=black,
  coltext=white,
  colframe=black,
  boxsep=5pt,
  arc=4pt}
\newtcolorbox{bonus}{
  colback=blue!15,
  colframe=blue!15,
  coltext=black!80,
  boxsep=5pt,
  arc=4pt}
\newtcolorbox{reflect}{
  colback=green!5,
  colframe=green!5,
  coltext=black!80,
  boxsep=5pt,
  arc=4pt}
\newtcolorbox{assessment}{
  colback=blue!5,
  colframe=blue!5,
  coltext=black!80,
  boxsep=5pt,
  arc=4pt}
\newtcolorbox{progress}{
  colback=purple!10,
  colframe=purple!10,
  coltext=black!80,
  boxsep=5pt,
  arc=4pt}
\newtcolorbox{video}{
  colback=yellow!5,
  colframe=yellow!5,
  coltext=black!80,
  boxsep=5pt,
  arc=4pt}
\newtcolorbox{caution}{
  colback=red!5,
  colframe=red!5,
  coltext=black!80,
  boxsep=5pt,
  arc=4pt}
\newtcolorbox{feedback}{
  colback=black!5,
  colframe=black!5,
  coltext=black!80,
  boxsep=5pt,
  arc=4pt}
\ifLuaTeX
  \usepackage{selnolig}  % disable illegal ligatures
\fi
\usepackage[]{natbib}
\bibliographystyle{apalike}
\IfFileExists{bookmark.sty}{\usepackage{bookmark}}{\usepackage{hyperref}}
\IfFileExists{xurl.sty}{\usepackage{xurl}}{} % add URL line breaks if available
\urlstyle{same}
\hypersetup{
  pdftitle={Philosophy for Life},
  hidelinks,
  pdfcreator={LaTeX via pandoc}}

\title{Philosophy for Life}
\author{}
\date{\vspace{-2.5em}}

\begin{document}
\maketitle

{
\setcounter{tocdepth}{1}
\tableofcontents
}
\hypertarget{welcome}{%
\chapter*{Welcome}\label{welcome}}
\addcontentsline{toc}{chapter}{Welcome}

This is the course book for Philosophy for Life. This book is divided into thematic units of study to help you engage with the materials. The course resources and learning activities are designed not only to help prepare you for the course assessments, but also to give you opportunities to practice various skills.

\begin{quote}
Below you will find information about how to navigate this book. Please read the full course syllabus located on the Course Home page in Moodle. It includes key information about the course schedule, assignments, and policies.
\end{quote}

\hypertarget{course-notes}{%
\section*{Course Notes}\label{course-notes}}
\addcontentsline{toc}{section}{Course Notes}

You should be reading this information in the context of a Trinity Western University course offered via Moodle. If this is not the case, then this may be an unauthorized reproduction of the course. Please contact \href{mailto:elearning@twu.ca}{\nolinkurl{elearning@twu.ca}} if you have concerns.

These notes will be your guide through the learning activities and assessment strategies necessary for you to succeed in the course, so it is important for you to engage to the best of your ability and take advantage of the resources available to you through Trinity Western University.

\begin{quote}
Assessment tasks are managed in other sections of the Moodle course, so be sure to familiarize yourself with those requirements and resources.
\end{quote}

\hypertarget{how-to-navigate-this-book}{%
\section*{How To Navigate This Book}\label{how-to-navigate-this-book}}
\addcontentsline{toc}{section}{How To Navigate This Book}

To move quickly to different portions of the book, click on the appropriate chapter or section in the table of contents on the left. The buttons at the top of the page allow you to show/hide the table of contents, search the book, change font settings, download a pdf or ebook copy of this book, or get hints on various sections of the book.

\includegraphics{assets/course-intro/menu.png}

The faint left and right arrows at the sides of each page (or bottom of the page if it's narrow enough) allow you to step to the next/previous section. Here's what they look like:

\includegraphics{assets/course-intro/left_arrow.png} \includegraphics{assets/course-intro/right_arrow.png}

You can also download an offline copy of this book in various formats, such as pdf or an ebook. If you are having any accessibility or navigation issues with this book, please reach out to your instructor or our online team at \href{mailto:elearning@twu.ca}{\nolinkurl{elearning@twu.ca}}.

\hypertarget{course-units}{%
\subsection*{Course Units}\label{course-units}}
\addcontentsline{toc}{subsection}{Course Units}

This course is organized into thematic units. Each unit of the course will provide you with the following information:

\begin{itemize}
\tightlist
\item
  A general overview of the key concepts that will be addressed during the unit.\\
\item
  Specific learning outcomes and topics for the unit.\\
\item
  Learning activities to help you engage with the concepts. These often include key readings, videos, and reflective prompts.\\
\item
  The Assessment section provides details on assignments you will need to complete throughout the course to demonstrate your understanding of the course learning outcomes.
\end{itemize}

\begin{quote}
Note that assessments, including assignments and discussion posts will be submitted in Moodle. See the Assessment section(s) in Moodle for full assignment details.
\end{quote}

\hypertarget{course-activities}{%
\subsection*{Course Activities}\label{course-activities}}
\addcontentsline{toc}{subsection}{Course Activities}

Below is some key information on features you will see throughout the course.

\begin{reflect}
\textbf{\emph{Learning Activity}}

This box will prompt you to engage in course concepts, often by viewing resources and reflecting on your experience and/or learning. Most learning activities are ungraded and are designed to help prepare you for the assessment in this course.
\end{reflect}

\begin{assessment}
\textbf{\emph{Assessment}}

This box will signify an assignment or discussion post you will submit in Moodle. Note that these demonstrate your understanding of the course learning outcomes. Be sure to review the grading rubrics for each assignment.
\end{assessment}

\begin{progress}
\textbf{\emph{Checking Your Learning}}

This box is for checking your understanding, to make sure you are ready for what follows. Ways to check your learning might include self-check quizzes or questions for discussion. These activities are not graded but are critical for you to be able to begin to develop evaluative judgement in this domain of knowledge.
\end{progress}

\begin{caution}
\textbf{\emph{Note}}

This box signifies key notes. It may also warn you of possible problems or pitfalls you may encounter!
\end{caution}

If you have any questions, do not hesitate to ask. We are here to help and be your guide on this journey.

\hypertarget{wisdom}{%
\chapter{Wisdom}\label{wisdom}}

This video briefly introduces the four topics in the first unit: knowledge, humility, friendship, and rhetoric. While we will learn some basic ideas about the nature of wisdom in the first topic on knowledge, most of this unit is devoted to introducing some of the skills of living a life of wisdom and how we can begin to practice these skills.

\begin{itemize}
\tightlist
\item
  \textbf{Wisdom Unit 1 Introduction Video} (6 min 10 sec)
\end{itemize}

\hypertarget{overview}{%
\section*{Overview}\label{overview}}
\addcontentsline{toc}{section}{Overview}

\hypertarget{topics}{%
\subsection*{Topics}\label{topics}}
\addcontentsline{toc}{subsection}{Topics}

This unit is divided into the following topics:

\begin{enumerate}
\def\labelenumi{\arabic{enumi}.}
\tightlist
\item
  Wisdom and Knowledge
\item
  Wisdom and Humility
\item
  Wisdom and Friendship
\item
  Wisdom and Rhetoric
\end{enumerate}

\hypertarget{learning-outcomes}{%
\subsection*{Learning Outcomes}\label{learning-outcomes}}
\addcontentsline{toc}{subsection}{Learning Outcomes}

When you've completed this unit, you will have learned how to:

\begin{itemize}
\tightlist
\item
  Identify and understand the role wisdom plays in knowledge, humiloity, friendship and rhetoric
\item
  Practice and apply the skills of wisdom to knowledge, humility, friendship and rhetoric
\item
  Apply these skills to the types of actions a wise person may take
\item
  Enhance your skills of wisdom and reason to live life within the context of justice and faith
\item
  Develop the skills to become a well-rounded person and citizen and live a wise and just life
\end{itemize}

\hypertarget{wisdom-and-knowledge}{%
\section*{1.1 Wisdom and Knowledge}\label{wisdom-and-knowledge}}
\addcontentsline{toc}{section}{1.1 Wisdom and Knowledge}

In this topic, we learn about the nature of wisdom and its relationship with knowledge. Wisdom is not just about knowledge. Wisdom is practical in that wisdom is required to live life well, to cope with and respond to the many challenges of life, and, if possible, avoid some of those challenges. While wisdom is practical, it does require at least some knowledge. For living life well, wisdom requires knowing the right beliefs and knowing how to apply these beliefs in the right way at the right time. How do we know which beliefs and actions are appropriate for living life well? Philosopher Robert Nozick argues that a wise person holds beliefs about many aspects of life, including beliefs about life's most important values, beliefs about achieving one's goals, and following the necessary steps to achieve these goals. Likewise, the right actions for living a fulfilling life are derived from applying the aforementioned beliefs. If a person knows about the most important values and goals in life, but never applies them to their behavior, then they are not considered wise. In other words, having knowledge does not entail wisdom. Nozick is also aware that applying the appropriate beliefs and actions is not sufficient to live life well. Sometimes bad things occur for no apparent reason, which may undermine our efforts. A wise person is aware of this problem, and while pursuing their goals, prepares for failure, learns how to avoid it, and how to respond to it if failure occurs.

\hypertarget{watch-and-reflect}{%
\subsection*{Watch and Reflect}\label{watch-and-reflect}}
\addcontentsline{toc}{subsection}{Watch and Reflect}

\begin{itemize}
\tightlist
\item
  Wisdom and Knowledge Unit 1 Topic 1 (11 min 15 sec)
\end{itemize}

\hypertarget{read-and-reflect}{%
\subsection*{Read and Reflect}\label{read-and-reflect}}
\addcontentsline{toc}{subsection}{Read and Reflect}

\begin{itemize}
\tightlist
\item
  Robert Nozick - \href{assets/u1/PHIL-100-Nozick-What-is-Wisdom.pdf}{\emph{What is Wisdom and Why Do Philosophers Love It So?}}
\end{itemize}

\hypertarget{wisdom-and-humility}{%
\section*{1.2 Wisdom and Humility}\label{wisdom-and-humility}}
\addcontentsline{toc}{section}{1.2 Wisdom and Humility}

In this topic, we focus on wisdom and humility. A wise person understands their limitations about what they know and do not know. In other words, a wise person is aware of their ignorance. The word ``ignorance'' in this case does not have a negative connotation, but rather describes a lack of knowledge about a particular thing.

In ``the Apology,'' Plato introduces the character Socrates, who is on trial for his life. During his defense, Socrates describes the account of his friend, Chaerephon, seeking answers from the Oracle of Delphi, wondering if there is any person wiser than Socrates. In reply to Chaerephon, the oracle says that no one is wiser than Socrates. Socrates is puzzled by the revelation about his wisdom since he considers himself to be unwise. Socrates then asks questions of different people in the city, such as politicians, poets, and artisans, in an attempt to solve the mystery about the oracle's revelation. Socrates discovers that while these individuals claim to know a great deal about issues such as the nature of beauty and what is good, he observes that they know very little about such issues. These people think they possess knowledge, when in fact they do not, and that is why they lack wisdom. Their way of thinking contrasts with Socrates, who thinks he lacks knowledge about these issues; more specifically, Socrates is aware of his ignorance about the nature of beauty and what is good, while the others are unaware of their own ignorance about these same issues. Because of this, Socrates is the wisest person.

One application of being aware of ignorance for seeking wisdom is to practice a humble attitude towards the complex issues in life, such as how to interpret the Bible, evolution, the nature of sexuality, the implications of social media for our lives, and was Donald Trump an effective president? These types of questions represent complex ideas that even some of the most specialized scholars cannot answer. A wise person demonstrates their humility towards these complex issues by asking thoughtful questions about the opposing view, knowing when to withhold judgment and opinion, and with a willingness to change their mind.

\hypertarget{watch-and-reflect-1}{%
\subsection*{Watch and Reflect}\label{watch-and-reflect-1}}
\addcontentsline{toc}{subsection}{Watch and Reflect}

\begin{itemize}
\tightlist
\item
  Wisdom and Humility Unit 1 Topic 2 (13 min 54 sec)
\end{itemize}

\hypertarget{read-and-reflect-1}{%
\subsection*{Read and Reflect}\label{read-and-reflect-1}}
\addcontentsline{toc}{subsection}{Read and Reflect}

\begin{itemize}
\tightlist
\item
  Plato - \emph{Apology} - \url{https://classics.mit.edu/Plato/apology.html}
\end{itemize}

\hypertarget{activity-practicing-humility---simulation-exercise-2}{%
\subsection*{Activity: Practicing Humility - Simulation Exercise 2}\label{activity-practicing-humility---simulation-exercise-2}}
\addcontentsline{toc}{subsection}{Activity: Practicing Humility - Simulation Exercise 2}

\begin{reflect}
The simulation exercises are meant to simulate a conversation and introduce you to the various topics in Unit 1. You are encouraged to not only review the simulation exercises but to and practice the applications of wisdom in your daily lives.
\end{reflect}

\hypertarget{watch-and-reflect-2}{%
\subsection*{Watch and Reflect}\label{watch-and-reflect-2}}
\addcontentsline{toc}{subsection}{Watch and Reflect}

\begin{reflect}
Please take a moment to watch this optional video with Dr.~Ian Church, What is intellectual humility? A doxastic account. Here is the first of three parts. If you have the time, you should watch them all. Please note that this video lasts for approximately 2:27 minutes.

\href{https://www.youtube.com/watch?v=8CZIkGEJYRY}{Watch: \emph{What is intellectual humility?}}
\end{reflect}

\hypertarget{wisdom-and-friendship}{%
\section*{1.3 Wisdom and Friendship}\label{wisdom-and-friendship}}
\addcontentsline{toc}{section}{1.3 Wisdom and Friendship}

This topic examines wisdom and friendship. A wise person is careful about how they view their friends. In doing so, it is best to learn about the various types of friendships. According to Aristotle, there are three types of friendships, friends of utility, pleasure, and virtue. Friends of utility are friends who are a means to an end, such as business partners, or some friends on a sports team. Friends of pleasure are friends with whom we share pleasurable activities. Friends of virtue are friends who exhibit good qualities and help us develop these qualities as well. Aristotle believes that the third type of friend is the best type of friend, and this type of friendship, while difficult to find and develop, is worth pursuing and fighting for. This is because friends of virtue care about us and our character development beyond the scope of utility and pleasure. They care about us even when things are difficult. But how do we find these friends and how can we be virtuous friends to other people? Thomas Aquinas provides some advice. Aquinas suggests that we should be careful about the bad qualities that make for a poor friend, such as envy and pride. These vices (or bad characteristics) undermine the pursuit of virtue and thus compromise any friendship that is built upon virtue. An envious person is someone who is upset about our own success. Think of a friend you have. Have they ever been upset or sad when good things happen to you? If the answer is yes, then this person may not be a virtuous friend. They may suffer from envy towards you. Or have you ever had a friend who jokes about you but can never receive the joke in return? This person may not be a virtuous friend because they suffer from arrogant pride. They think too highly of themselves. Furthermore, we should be cautious about whether we are exhibiting these defects as well. Do we treat our friends as a means to a utilitarian end? Do we make friends only for pleasure? Do we feel upset about our friends success? Do we exhibit a sense of arrogant pride about ourselves? This video introduces these ideas from Aristotle and Aquinas to help better equip us with the skills for developing healthy friendship for living well.

\hypertarget{watch-and-reflect-3}{%
\subsection*{Watch and Reflect}\label{watch-and-reflect-3}}
\addcontentsline{toc}{subsection}{Watch and Reflect}

\begin{itemize}
\tightlist
\item
  Wisdom and Friendship Unit 1 Topic 3 (24 min 18 sec)
\end{itemize}

\hypertarget{read-and-reflect-2}{%
\subsection*{Read and Reflect}\label{read-and-reflect-2}}
\addcontentsline{toc}{subsection}{Read and Reflect}

\begin{itemize}
\tightlist
\item
  Aristotle - \href{assets/u1/PHIL-100-Aristotle-NE-VIII-IX-X.pdf}{\emph{Nicomachean Ethics}, Books VIII, IX and X}
\end{itemize}

\hypertarget{watch-and-reflect-4}{%
\subsection*{Watch and Reflect}\label{watch-and-reflect-4}}
\addcontentsline{toc}{subsection}{Watch and Reflect}

\begin{reflect}
Please take a moment to view this optional video. In this video, the Honorable Richard Holloway, who is a writer, broadcaster, and cleric from Scotland, seeks to initiate a debate around the more profound spiritual questions of life. The duration of this video is 4:45 minutes.

\href{https://www.youtube.com/watch?v=LERi_Xjfio4}{Watch: \emph{Why doubt is a good thing}}
\end{reflect}

The aim of the following learning activities is to explore and learn about these two applications of knowledge: the practice of doubt and the acquisition of valuable knowledge over entertainment.

\hypertarget{activity-learning-to-doubt---simulation-exercise-1}{%
\subsection*{Activity : Learning to doubt - Simulation Exercise 1}\label{activity-learning-to-doubt---simulation-exercise-1}}
\addcontentsline{toc}{subsection}{Activity : Learning to doubt - Simulation Exercise 1}

\begin{reflect}
The simulation exercises are meant to simulate a conversation and introduce students to the various topics in Unit 1. Students are encouraged to not only review the simulation exercises but to practice the applications of wisdom in their daily lives.
\end{reflect}

\hypertarget{activity-take-notes-for-your-life-journal}{%
\subsection*{Activity: Take notes for your Life Journal}\label{activity-take-notes-for-your-life-journal}}
\addcontentsline{toc}{subsection}{Activity: Take notes for your Life Journal}

\begin{reflect}
Answering the following questions is optional. Answering these questions and saving them as a document will allow you to take notes on the readings in preparation for the Life Journal assignment due at the end of Unit 1.
\end{reflect}

\hypertarget{wisdom-and-rhetoric}{%
\section*{1.4 Wisdom and Rhetoric}\label{wisdom-and-rhetoric}}
\addcontentsline{toc}{section}{1.4 Wisdom and Rhetoric}

This topic examines wisdom and rhetoric. Arguments with other people, such as family, friends, or colleagues, are part of everyday life. A wise person knows how to argue without fighting. According to Jay Heinrichs, a necessary skill for arguing productively is knowing the appropriate goal you wish to bring about. Are you trying to influence someone's decision? Is that feasible? Are you hoping to feel like a winner? Is that an appropriate goal? It may be your goal to change the mood of your opponent, or simply to inform your opponent of a different point of view. Would you like to end the argument quickly? After the argument, would you like to maintain the relationship with the person? Should you pursue either of these goals? How will you go about accomplishing your goal? In what way are you prepared to compromise? Are you prepared to experience defeat from your opponent's point of view in order to achieve your goal? What is the purpose of the argument and how should it be achieved?

\hypertarget{watch-and-reflect-5}{%
\subsection*{Watch and Reflect}\label{watch-and-reflect-5}}
\addcontentsline{toc}{subsection}{Watch and Reflect}

\begin{itemize}
\tightlist
\item
  Wisdom and Rhetoric Unit 1 Topic 4 (12 min 19 sec)
\end{itemize}

\hypertarget{read-and-reflect-3}{%
\subsection*{Read and Reflect}\label{read-and-reflect-3}}
\addcontentsline{toc}{subsection}{Read and Reflect}

\begin{itemize}
\tightlist
\item
  Jay Heinrichs - \href{assets/u1/PHIL-100-Heinrichs-Thank-You-for-Arguing.pdf}{\emph{Thank You for Arguing}}
\end{itemize}

\hypertarget{activity-take-notes-for-your-life-journal-1}{%
\subsection*{Activity: Take notes for your Life Journal}\label{activity-take-notes-for-your-life-journal-1}}
\addcontentsline{toc}{subsection}{Activity: Take notes for your Life Journal}

\begin{reflect}
Answering the following questions is optional. Answering these questions and saving them as a document will allow you to take notes on the readings in preparation for the Life Journal assignment due at the end of Unit 1.
\end{reflect}

\hypertarget{watch-and-reflect-6}{%
\subsection*{Watch and Reflect}\label{watch-and-reflect-6}}
\addcontentsline{toc}{subsection}{Watch and Reflect}

\begin{reflect}
Take a moment to watch this optional video. Here, Camille A. Langston, a specialist in rhetoric and nineteenth-century women, explains the basics of deliberative rhetoric and shares tips for appealing to the ethos, logos, and pathos of your audience. This video is 4:29 minutes long.

\href{https://www.youtube.com/watch?v=3klMM9BkW5o}{Watch: \emph{How to use rhetoric to get what you want}}
\end{reflect}

\hypertarget{required-note-taking-learning-activity}{%
\section*{Required Note Taking Learning Activity}\label{required-note-taking-learning-activity}}
\addcontentsline{toc}{section}{Required Note Taking Learning Activity}

\begin{reflect}
\begin{itemize}
\tightlist
\item
  Answering these questions will help you reflect on the video topic content and prepare you for your Unit 1 Reflection Assignment.
\item
  Click the Download button after you've completed all the questions for this unit and save the answers to your computer.
\item
  You will need the downloaded answers to these questions as you will be required to submit them as part of your Final Course Reflection Essay/Video Assignment.
\end{itemize}
\end{reflect}

\hypertarget{faith}{%
\chapter{Faith}\label{faith}}

\hypertarget{overview-1}{%
\section*{Overview}\label{overview-1}}
\addcontentsline{toc}{section}{Overview}

This unit is divided into three lectures.

\begin{enumerate}
\def\labelenumi{\arabic{enumi}.}
\tightlist
\item
  Wisdom and Humility Lecture 1 discusses an allegory.
\item
  Wisdom and Humility Lecture 2 is a section from Plato.
\item
  Wisdom and Charity Lecture discusses the importance of first defending your opponent's view before raising objections.
\end{enumerate}

\hypertarget{learning-outcomes-1}{%
\subsection*{Learning Outcomes}\label{learning-outcomes-1}}
\addcontentsline{toc}{subsection}{Learning Outcomes}

When you've completed this unit, you will have learned how to:

\begin{itemize}
\tightlist
\item
  Apply wisdom to practice intellectual humility with knowledge
\item
  Interpret Plato's ``The Apology'' from a wise person's perspective
\item
  Develop and apply wisdom to strategies for practicing charity towards your opponents
\end{itemize}

\hypertarget{wisdom-and-humility-lecture-1}{%
\section*{2.1 Wisdom and Humility Lecture 1}\label{wisdom-and-humility-lecture-1}}
\addcontentsline{toc}{section}{2.1 Wisdom and Humility Lecture 1}

This topic examines the first part about wisdom and humility. If wisdom includes learning and practicing the right types of knowledge and avoiding the wrong types of knowledge, then one ought to practice humility. One must find a way to accept their limitations about what they know. Indeed, the more one learns, often the more they realize how much they do not know. Thus, knowledge often extends outwards, growing more and more as one learns the topic. A wise person practices the skill of not caring about how much they think they know. A wise person practices the skill of not caring when they are wrong. Indeed, they recognize their error as the opportunity to grow in their knowledge. In the following video, we explore the allegory of the Knowledge Rooms to help us imagine the challenges of obtaining knowledge and the importance of maintaining a humble disposition.

Read

\begin{itemize}
\tightlist
\item
  Plato's \emph{Apology} - \url{https://classics.mit.edu/Plato/apology.html}
\end{itemize}

Watch

\begin{itemize}
\tightlist
\item
  Wisdom and Humility - Unit 2 Lecture 2.1 video (13 min 12 sec)
\item
  This video introduces wisdom and humility. Intellectual humility is valuable for keeping our so-called knowledge in check. A wise person is aware of their own ignorance. This idea is made clearer with the allegory of the knowledge rooms. The term allegory means a kind of story to represent a deeper point.
\end{itemize}

\hypertarget{wisdom-and-humility-lecture-2}{%
\section*{2.2 Wisdom and Humility Lecture 2}\label{wisdom-and-humility-lecture-2}}
\addcontentsline{toc}{section}{2.2 Wisdom and Humility Lecture 2}

This topic examines the second part about wisdom and humility. Plato believed that a wise person understands the limitations of what they know and do not know. In other words, a wise person is aware of their ignorance. The word ``ignorance'' in this case does not have a negative connotation, but rather describes a lack of knowledge about a particular thing.

In ``the Apology of Socrates,'' Plato introduces the character Socrates, who is on trial for his life. During his defense, Socrates describes the account of his friend, Chaerephon, seeking answers from the Oracle of Delphi, wondering if there is any person wiser than Socrates. In reply to Chaerephon, the oracle says that no one is wiser than Socrates. Socrates is puzzled by the revelation about his wisdom since he considers himself to be unwise. Socrates then asks questions of different people in the city, such as politicians, poets, and artisans, in an attempt to solve the mystery about the oracle's revelation. Socrates discovers that while these individuals claim to know a great deal about issues such as the nature of beauty and what is good, he observes that they know very little about such issues. These people think they possess knowledge, when in fact they do not, and that is why they lack wisdom. Their way of thinking contrasts with Socrates, who thinks he lacks knowledge about these issues; more specifically, Socrates is aware of his ignorance about the nature of beauty and what is good, while the others are unaware of their own ignorance about these same issues. Because of this, Socrates is the wisest person.

One application of being aware of ignorance for seeking wisdom is to practice a humble attitude towards the complex issues in life, such as how to interpret the Bible, evolution, the nature of sexuality, the implications of social media for our lives, and was Donald Trump an effective president? These types of questions represent complex ideas that even some of the most specialized scholars cannot answer. A wise person demonstrates their humility towards these complex issues by asking thoughtful questions about the opposing view, knowing when to withhold judgment and opinion, and with a willingness to change their mind. In the following video, we explore Plato's description of Socrates when confronted by his accusers and how Socrates seems to exhibit humility.

Watch

\begin{itemize}
\tightlist
\item
  Plato's \emph{The Apology of Socrates} - Unit 2 Lecture 2.2 video (12 min 57 sec)
\item
  This video expands on wisdom and humility. This includes a short section of Plato's ``The Apology of Socrates,'' and reinforces the idea that a wise person is aware of their own ignorance. I wise person knows about their own skills and competency, but is also aware of the many things they do not know, such as their ideas about politics, sexuality, religion, and science.
\end{itemize}

\hypertarget{wisdom-and-charity}{%
\section*{2.3 Wisdom and Charity}\label{wisdom-and-charity}}
\addcontentsline{toc}{section}{2.3 Wisdom and Charity}

This topic examines wisdom and charity. The term ``charity'' in this case refers to a principle of treating a person's argument and/or opinion in its strongest form. Charity involves more than humility. It requires willpower. One must willfully, and temporarily, embrace the opponent's view, assist with constructing that view in its best possible light, and then express the view better than the opponent. The act of embracing an opposing argument is challenging for many reasons, but one reason is that humans are not built to act with charity toward those with whom they disagree, especially when those convictions are strong. Put in a different way, humans struggle to, on the one hand, affirm a viewpoint in which they hold strong convictions, and then, on the other hand, compartmentalize that view away while attempting to charitably evaluate objections to that conviction.

A wise person practices charity because of their humility. They recognize the value of charity because they recognize the limitations of their knowledge. That is, perhaps the opponent's view is correct. Or perhaps the spirit of the opponent's view is correct, even if the opponent cannot properly state this view. Charity requires that we help this opponent first to be sure their view is represented accurately before raising counterpoints. In the following video, we examine this idea of charity and some of the problems with practicing charity.

Watch

\begin{itemize}
\tightlist
\item
  Wisdom and Charity - Unit 2 Lecture 2.3 video (13 min 36 sec)
\item
  This video introduces wisdom and charity. The term charity in this case applies to intellectual charity, which means presenting your opponents view in its strongest light before raising objections. A wise person does not take advantage of their opponent, such as their inability to communicate well, or the mistakes they make. Rather, a charitable person helps their opponent build their best case and then raise objections afterwards.
\end{itemize}

\hypertarget{required-note-taking-learning-activity-1}{%
\section*{Required Note Taking Learning Activity}\label{required-note-taking-learning-activity-1}}
\addcontentsline{toc}{section}{Required Note Taking Learning Activity}

\begin{itemize}
\tightlist
\item
  Answering these questions will help you reflect on the video lecture content and prepare you for your Unit 2 Reflection Assignment.
\item
  Click the Download button after you've completed all the questions for this unit and save the answers to your computer.
\item
  You will be required to submit these downloaded notes as part of your Final Course Reflection Essay/Video Assignment.
\end{itemize}

\hypertarget{reason}{%
\chapter{Reason}\label{reason}}

\hypertarget{overview-2}{%
\section*{Overview}\label{overview-2}}
\addcontentsline{toc}{section}{Overview}

This unit is divided into four lectures.

\begin{enumerate}
\def\labelenumi{\arabic{enumi}.}
\tightlist
\item
  Wisdom and Friendship Lecture 1 on Aristotle's view of friendship.
\item
  Wisdom and Friendship Lecture 2 on Aquinas' views of pride and envy and how they affect friendship.
\item
  Wisdom and Rhetoric Lecture 1 focusses on the importance of rhetoric.
\item
  Wisdom and Rhetoric Lecture 2 focusses on some skills for practicing rhetoric.
\end{enumerate}

\hypertarget{learning-outcomes-2}{%
\subsection*{Learning Outcomes}\label{learning-outcomes-2}}
\addcontentsline{toc}{subsection}{Learning Outcomes}

When you've completed this unit, you will have learned how to:

\begin{itemize}
\tightlist
\item
  Apply wisdom to identify what forms a good friendship and a bad friendship
\item
  Identify and apply the importance of virtue in friendship
\item
  Apply wisdom to identify and use some basic skills of rhetoric (or persuasion)
\end{itemize}

\hypertarget{wisdom-and-friendship-lecture-1}{%
\section*{3.1 Wisdom and Friendship Lecture 1}\label{wisdom-and-friendship-lecture-1}}
\addcontentsline{toc}{section}{3.1 Wisdom and Friendship Lecture 1}

This topic examines the first part about wisdom and friendship. A wise person is careful about how they view their friends. In doing so, it is best to learn about the various types of friendships. According to Aristotle, there are several types of friendships, such as friends of utility, pleasure, and virtue. Friends of utility are friends who are a means to an end, such as business partners, or some friends on a sports team. Friends of pleasure are friends with whom we share pleasurable activities. Friends of virtue are friends who exhibit good qualities and help us develop these qualities as well. Aristotle believes that the third type of friend is the best friend, and this type of friendship, while difficult to find and develop, is worth pursuing and fighting for. In the following video, we briefly outline Aristote's view of friendship and then figure how a wise person should think about friendship in these ways.

Watch

\begin{itemize}
\tightlist
\item
  Wisdom and Friendship - Unit 3 Lecture 3.1 video (14 min 14 sec)
\item
  This video introduces Aristotle's view of friendship. A wise person is aware of the different types of friendships and the importance of these types. This video helps us understand what makes a good friend and how we can be good friends to other people.
\end{itemize}

\hypertarget{wisdom-and-friendship-lecture-2}{%
\section*{3.2 Wisdom and Friendship Lecture 2}\label{wisdom-and-friendship-lecture-2}}
\addcontentsline{toc}{section}{3.2 Wisdom and Friendship Lecture 2}

This topic examines the second part about wisdom and friendship. For Aristotle, friends of virtue are the best of friends because they care about us and our character development beyond the scope of utility and pleasure. They care about us even when things are difficult. But how do we find these friends and how can we be virtuous friends to other people? Thomas Aquinas provides some advice. Aquinas suggests that we should be careful about the bad qualities that make for a poor friend, such as envy and pride. These vices (or bad characteristics) undermine the pursuit of virtue and thus compromise any friendship that is built upon virtue.

An envious person is someone who is upset about our own success. Think of a friend you have. Have they ever been upset or sad when good things happen to you? If the answer is yes, then this person may not be a virtuous friend. They may suffer from envy towards you. Or have you ever had a friend who jokes about you but can never receive the joke back to them in turn? This person may not be a virtuous friend because they suffer from arrogant pride. They think too highly of themselves. Furthermore, we should be cautious about whether we are exhibiting these defects as well. Do we feel upset about our friends success? Do we exhibit a sense of arrogant pride about ourselves? These vices undermine virtuous friendships. In the following video, we briefly outline Aquinas' views about envy and pride and note how these vices undermine our friendships.

Watch

\begin{itemize}
\tightlist
\item
  Wisdom and Friendship - Unit 3 Lecture 3.2 video (10 min 19 sec)
\item
  This video introduces Aquinas' view of pride and envy. These vices are damaging to friendships because they undermine virtue. This video helps us understand how to identify bad friends and how we can be aware of when we are being a bad friend to others.
\end{itemize}

\hypertarget{wisdom-and-rhetoric-lecture-1}{%
\section*{3.3 Wisdom and Rhetoric Lecture 1}\label{wisdom-and-rhetoric-lecture-1}}
\addcontentsline{toc}{section}{3.3 Wisdom and Rhetoric Lecture 1}

This topic examines the first part about wisdom and rhetoric. Rhetoric refers to the art of persuasion or influence. Previously we discussed how wisdom requires an awareness and application of our actions, knowledge, humility, charity, and friendship making. The art of persuasion applies to all of these skills. Rhetoric can help us get along in life in these various domains, particularly when dealing with other people.

But many people throughout history have neglected this art of rhetoric. Indeed, Plato thought rhetoric was for fools. It was the art of manipulation. Rhetoric was for the sophists and those who only wished to confuse the masses, aiming these people away from the truth. However, his student, Aristotle, rejected Plato's view. For Aristotle, rhetoric was a tool that, while available for manipulation, was equally valuable for communicating the truth. Rhetoric was the only weapon suitable to combat such manipulation. Thus, it is better to hone the skills of rhetoric to combat the sophists. In the following video, we analyze Plato's objections to rhetoric and Aristotle's reasons for rejecting Plato's views and embracing the art of rhetoric.

Read

\begin{itemize}
\tightlist
\item
  Jay Heinrichs - \emph{Thank You for Arguing}
\end{itemize}

Watch

\begin{itemize}
\tightlist
\item
  Wisdom and Rhetoric - Unit 3 Lecture 3.3 video (8 min 01 sec)
\item
  This video introduces wisdom and rhetoric. Rhetoric is the art of persuasion. This is a valuable skill to practice in life because we often engage in rhetoric with other people, family, friends, supervisors, strangers, and even enemies. This video introduces Aristotle's defense of using rhetoric against Plato's criticisms of the practice.
\end{itemize}

\hypertarget{wisdom-and-rhetoric-lecture-2}{%
\section*{3.4 Wisdom and Rhetoric Lecture 2}\label{wisdom-and-rhetoric-lecture-2}}
\addcontentsline{toc}{section}{3.4 Wisdom and Rhetoric Lecture 2}

This topic examines the second part about wisdom and rhetoric. Arguments with other people, such as family, friends, or colleagues, are part of everyday life. A wise person knows how to argue without fighting. According to Jay Heinrichs, a necessary skill for arguing productively is knowing the appropriate goal you wish to bring about. Are you trying to influence someone's decision? Is that feasible? Are you hoping to feel like a winner? Is that an appropriate goal? It may be your goal to change the mood of your opponent, or simply to inform your opponent of a different point of view. Would you like to end the argument quickly? After the argument, would you like to maintain the relationship with the person? Should you pursue either of these goals? How will you go about accomplishing your goal? In what way are you prepared to compromise? Are you prepared to experience defeat from your opponent's point of view in order to achieve your goal? What is the purpose of the argument and how should it be achieved?

Exercising the previous skills of wisdom when arguing with other people is helpful for practicing rhetoric. Wise actions, for example, require setting reasonable goals and avoiding the pitfalls that may undermine these goals. Similarly, when applying rhetoric in an argument with someone else, focus on the goal you wish to bring about and be willing to concede as much as possible to get there. Furthermore, remember that wisdom requires knowledge and humility. When these skills are applied to rhetoric, demonstrate the ability to challenge your own views to show your intention to understand the opponent's position. And in everything, act with charity. When arguing with others, assist them with their arguments, and, if possible, suggest ways in which their arguments can be strengthened. Demonstrate your comfort with the exchange by temporarily embracing their viewpoint. Charity of this kind encourages your opponent to act in a similar manner, and, in doing so, may reduce the defensiveness of each person. Depending on your rhetorical goal when arguing with another person, the tactics above may be useful. The following video explores some of these simple rhetorical strategies to help us better get along with other people.

Watch

\begin{itemize}
\tightlist
\item
  Wisdom and Rhetoric - Unit 3 Lecture 3.4 video (11 min 13 sec)
\item
  This video expands on the study of rhetoric by discussing some strategies for applying rhetoric. One strategy is to set goals of the encounter and focus on what you must do to achieve these goals. Perhaps you want the other person to understand your view, or perhaps you want to end the encounter with the person as quickly as possible without making things worse. Practicing rhetoric can help us in these ways.
\end{itemize}

\hypertarget{required-note-taking-learning-activity-2}{%
\section*{Required Note Taking Learning Activity}\label{required-note-taking-learning-activity-2}}
\addcontentsline{toc}{section}{Required Note Taking Learning Activity}

\begin{itemize}
\tightlist
\item
  Answering these questions will help you reflect on the video lecture content and prepare you for your Unit 3 Reflection Assignment.
\item
  Click the Download button after you've completed all the questions for this unit and save the answers to your computer.
\item
  You will be required to submit these downloaded notes as part of your Final Course Reflection Essay/Video Assignment.
\end{itemize}

\hypertarget{assessment}{%
\chapter*{Assessment}\label{assessment}}
\addcontentsline{toc}{chapter}{Assessment}

The following assignments are opportunities for learners to demonstrate their understanding of the course outcomes. Please confirm assignment details with your instructor, referring to the course syllabus.

Note that Assignment dropboxes are located in Moodle. Also refer to the Course Schedule in Moodle for the specific due dates.

\hypertarget{assignment}{%
\section*{Assignment:}\label{assignment}}
\addcontentsline{toc}{section}{Assignment:}

\hypertarget{grading-criteria}{%
\subsection*{Grading Criteria}\label{grading-criteria}}
\addcontentsline{toc}{subsection}{Grading Criteria}

See the following rubric that explains how your assignment will be evaluated. Also available as a \href{assets/assessment/Identity-as-a-Teacher-RUBRIC.pdf}{pdf}

\textbf{APA/WRITING}

\textbf{Unsatisfactory:} Paper does not model language and conventions used in scholarly literature. Writing is not well-organized. Several errors in grammar or composition. Sources are not cited. APA citations are not appropriately formatted.

\textbf{Developing:} Paper partially models language and conventions used in scholarly literature. Writing is somewhat well organized and includes some errors in grammar or composition. Not all sources cited. APA citations are generally formatted correctly, with several errors.

\textbf{Proficient:} \emph{Paper consistently models language and conventions used in scholarly literature. Writing is well-organized and includes few (if any) errors in grammar or composition. All resources are appropriately cited (including in-text citations and bibliography information). Few (if any) errors in APA citations.}

\textbf{Exemplary:} Paper is an exemplar of language and conventions used in scholarly literature. Writing is well-organized and free of errors in grammar or composition. All resources are appropriately cited. No errors in APA format.

\hypertarget{statement-of-teaching-identity}{%
\subsubsection*{STATEMENT OF TEACHING IDENTITY}\label{statement-of-teaching-identity}}
\addcontentsline{toc}{subsubsection}{STATEMENT OF TEACHING IDENTITY}

\textbf{Unsatisfactory:} Does not provide a statement about identity as a teacher/facilitator

\textbf{Developing:} Provides an unclear statement about identity as a teacher/facilitator.

\textbf{Proficient:} \emph{Provides a clear, concise, and powerful statement about identity as a teacher/facilitator.}

\textbf{Exemplary:} Provides a clear, concise, and powerful statement about identity as a teacher/facilitator. Statement incorporates theory or research from course materials.

\hypertarget{developing-a-cohesive-and-logical-academic-argument}{%
\subsubsection*{DEVELOPING A COHESIVE AND LOGICAL ACADEMIC ARGUMENT}\label{developing-a-cohesive-and-logical-academic-argument}}
\addcontentsline{toc}{subsubsection}{DEVELOPING A COHESIVE AND LOGICAL ACADEMIC ARGUMENT}

\textbf{Unsatisfactory:} Does not make a focused, cohesive, or logical academic argument. Paper is confusing, and is missing an introduction, body, or conclusion. Transitions between sections and ideas are missing.

\textbf{Developing:} Makes an academic argument that is only partially focused, cohesive and logical. Paper is generally organized, but is missing an introduction, body, or conclusion. Transitions between sections and ideas are unclear.

\textbf{Proficient:} \emph{Makes a focused, cohesive, logical academic argument. Paper is effectively organized and includes an introduction, body, and conclusion. Transitions between sections and ideas are clear.}

\textbf{Exemplary:} Makes a focused, cohesive, logical and compelling academic argument. Paper is effectively organized and includes an introduction, body, and conclusion. Transitions between sections and ideas are clear, and build on each other.

\hypertarget{analysis-of-identity-as-a-teacher}{%
\subsubsection*{ANALYSIS OF IDENTITY AS A TEACHER}\label{analysis-of-identity-as-a-teacher}}
\addcontentsline{toc}{subsubsection}{ANALYSIS OF IDENTITY AS A TEACHER}

\textbf{Unsatisfactory:} Does not include three important aspects of identity as a teacher/facilitator. Does not include an analysis.

\textbf{Developing:} Lists but does not discuss three important aspects of identity as a teacher/facilitator. Includes a partial analysis.

\textbf{Proficient:} \emph{Includes a detailed discussion of three important aspects of identity as a teacher/facilitator. Includes thoughtful analysis of each of the three elements.}

\textbf{Exemplary:} Includes a detailed discussion of three important aspects of identity as a teacher/facilitator. Includes a thoughtful analysis, integrating scholarly literature to support analysis and furthering scholarly thinking related to teacher identity.

\hypertarget{scholarly-integration}{%
\subsubsection*{SCHOLARLY INTEGRATION}\label{scholarly-integration}}
\addcontentsline{toc}{subsubsection}{SCHOLARLY INTEGRATION}

\textbf{Unsatisfactory:} Does not integrate references to support claims and assertions made in the paper.

\textbf{Developing:} Integrates references to support some of the claims and assertions made in the paper.

\textbf{Proficient:} \emph{Integrates references to support claims and assertions made in the paper.}

\textbf{Exemplary:} Integrates references to support claims and assertions made in the paper, effectively synthesizing different perspectives and research results from scholarly sources.

\begin{longtable}[]{@{}
  >{\raggedright\arraybackslash}p{(\columnwidth - 8\tabcolsep) * \real{0.2000}}
  >{\raggedright\arraybackslash}p{(\columnwidth - 8\tabcolsep) * \real{0.2000}}
  >{\raggedright\arraybackslash}p{(\columnwidth - 8\tabcolsep) * \real{0.2000}}
  >{\raggedright\arraybackslash}p{(\columnwidth - 8\tabcolsep) * \real{0.2000}}
  >{\raggedright\arraybackslash}p{(\columnwidth - 8\tabcolsep) * \real{0.2000}}@{}}
\toprule\noalign{}
\begin{minipage}[b]{\linewidth}\raggedright
\textbf{TOTAL}
\end{minipage} & \begin{minipage}[b]{\linewidth}\raggedright
\textbf{0 = 0\% (F)}
\end{minipage} & \begin{minipage}[b]{\linewidth}\raggedright
\textbf{10 = 50\% (C)}
\end{minipage} & \begin{minipage}[b]{\linewidth}\raggedright
\textbf{15 = 75 (B)}
\end{minipage} & \begin{minipage}[b]{\linewidth}\raggedright
\textbf{20 = 100\% (A+)}
\end{minipage} \\
\midrule\noalign{}
\endhead
\bottomrule\noalign{}
\endlastfoot
\end{longtable}

\begin{center}\rule{0.5\linewidth}{0.5pt}\end{center}

\hypertarget{assignment-company-website-analysis}{%
\section*{Assignment: Company Website Analysis}\label{assignment-company-website-analysis}}
\addcontentsline{toc}{section}{Assignment: Company Website Analysis}

\begin{assessment}
Investigate the Human Resources or Faculty Development portion of a
company's website, a higher education institution or adult learning
facility, preferably one with which you are familiar. Focus on the
faculty or employee development part of the website. In this assignment,
you will apply the theory of teaching in/for/with depth by analyzing the
learning culture of an organization.

In a 4-5 page APA formatted paper, analyze the website by responding to
the following questions in your report:

\begin{enumerate}
\def\labelenumi{\arabic{enumi}.}
\tightlist
\item
  What can you infer about the company's learning culture?\\
\item
  From what is visible on the public website, would you say it is an
  authentic learning community? Why or why not? Discuss whether the
  website reflects aspects of one or more of the learning community
  models explored in previous lessons.\\
\item
  Do you see evidence that interconnectedness and integrity are valued?
  Explain.\\
\item
  What traits and skills seem to be valued in employees?\\
\item
  How does the company develop skills in its employees (e.g., workshops,
  seminars, mentoring)? Are the methods based on the principles of
  andragogy? (see Smith YouTube video). What specific adult learning
  strategies do you see reflected in the development/training
  opportunities for employees?
\end{enumerate}

Your paper should be 4-5 pages and should incorporate references to at
least five scholarly sources you have studied in this course, or other
scholarly sources you have identified.

The paper should include:

\begin{enumerate}
\def\labelenumi{\arabic{enumi}.}
\tightlist
\item
  Introduction\\
\item
  Analysis (responding to the prompts)\\
\item
  Conclusion\\
\item
  Reference List
\end{enumerate}
\end{assessment}

\hypertarget{company-website-analysis-rubric}{%
\subsection*{Company Website Analysis Rubric}\label{company-website-analysis-rubric}}
\addcontentsline{toc}{subsection}{Company Website Analysis Rubric}

See the following rubric that explains how your assignment will be evaluated. Also available as a \href{assets/assessment/Company-Website-Analysis-RUBRIC.pdf}{pdf}

\hypertarget{apa-formatting}{%
\subsubsection*{APA Formatting}\label{apa-formatting}}
\addcontentsline{toc}{subsubsection}{APA Formatting}

\textbf{Unsatisfactory:} Paper does not model language and conventions used in scholarly literature.
Writing is not well-organized. Several errors in grammar or composition. Sources
are not cited. APA citations are not appropriately formatted.

\textbf{Developing:} Paper partially models language and conventions used in scholarly literature.
Writing is somewhat well organized and includes some errors in grammar or
composition. Not all sources cited. APA citations are generally formatted
correctly, with several errors.

\textbf{Proficient:} \emph{Paper consistently models language and conventions used in scholarly
literature. Writing is well-organized and includes few (if any) errors in
grammar or composition. All resources are appropriately cited (including in-text
citations and bibliography information). Few (if any) errors in APA citations.}

\textbf{Exemplary:} Paper is an exemplar of language and conventions used in scholarly literature.
Writing is well-organized and free of errors in grammar or composition. All
resources are appropriately cited. No errors in APA format.

\hypertarget{developing-a-cohesive-and-logical-academic-argument-1}{%
\subsubsection*{DEVELOPING a COHESIVE and LOGICAL ACADEMIC ARGUMENT}\label{developing-a-cohesive-and-logical-academic-argument-1}}
\addcontentsline{toc}{subsubsection}{DEVELOPING a COHESIVE and LOGICAL ACADEMIC ARGUMENT}

\textbf{Unsatisfactory:} Does not make a focused, cohesive, or logical academic argument. Paper is
confusing, and is missing an introduction, body, or conclusion. Transitions
between sections and ideas are missing.

\textbf{Developing:} Makes an academic argument that is only partially focused, cohesive and logical.
Paper is generally organized, but is missing an introduction, body, or
conclusion. Transitions between sections and ideas are unclear.

\textbf{Proficient:} \emph{Makes a focused, cohesive, logical academic argument. Paper is effectively
organized and includes an introduction, body, and conclusion. Transitions
between sections and ideas are clear.}

\textbf{Exemplary:} Makes a focused, cohesive, logical and compelling academic argument. Paper is
effectively organized and includes an introduction, body, and conclusion.
Transitions between sections and ideas are clear and build on each other.

\hypertarget{analysis-of-learning-culture}{%
\subsubsection*{ANALYSIS of LEARNING CULTURE}\label{analysis-of-learning-culture}}
\addcontentsline{toc}{subsubsection}{ANALYSIS of LEARNING CULTURE}

\textbf{Unsatisfactory:} Does not include an analysis of the company learning culture, and no evaluation
of the authenticity of the learning community.

\textbf{Developing:} Includes a partial analysis of the company learning culture, including a limited
evaluation of the authenticity of the learning community.

\textbf{Proficient:} \emph{Includes a detailed analysis of the company learning culture, including an
evaluation of the authenticity of the learning community.}

\textbf{Exemplary:} Includes a detailed analysis of the company learning culture, including an
evaluation of the authenticity of the learning community. Includes a thoughtful
analysis, integrating scholarly literature to support analysis and furthering
scholarly thinking related to teacher identity.

\hypertarget{evaluation-of-interconnectedness-and-integrity}{%
\subsubsection*{EVALUATION of INTERCONNECTEDNESS and INTEGRITY}\label{evaluation-of-interconnectedness-and-integrity}}
\addcontentsline{toc}{subsubsection}{EVALUATION of INTERCONNECTEDNESS and INTEGRITY}

\textbf{Unsatisfactory:} Does not include an evaluation of evidence of interconnectedness and integrity
on the company website. Does not integrate scholarly sources in the evaluation.

\textbf{Developing:} Includes a partial evaluation of evidence of interconnectedness and integrity on
the company website. Evaluation includes only limited reference to scholarly
sources.

\textbf{Proficient:} \emph{Includes a detailed evaluation of evidence of interconnectedness and integrity
on the company website. Evaluation integrates scholarly sources.}

\textbf{Exemplary:} Includes a detailed evaluation of evidence of interconnectedness and integrity
on the company website. Includes recommendations for ways in which to integrate
interconnectedness and integrity into employee development.

\hypertarget{analysis-of-adult-learning-strategies}{%
\subsubsection*{ANALYSIS of ADULT LEARNING STRATEGIES}\label{analysis-of-adult-learning-strategies}}
\addcontentsline{toc}{subsubsection}{ANALYSIS of ADULT LEARNING STRATEGIES}

\textbf{Unsatisfactory:} Does not include a detailed analysis of valued skills and evidence of adult
learning theory in employee development. Does not integrate scholarly sources.

\textbf{Developing:} Includes a partial analysis of valued skills and evidence of adult learning
theory in employee development. Analysis integrates few, if any, scholarly
sources.

\textbf{Proficient:} \emph{Includes a detailed analysis of valued skills and evidence of adult learning
theory in employee development. Analysis integrates scholarly sources.}

\textbf{Exemplary:} Includes a detailed analysis of valued skills and evidence of adult learning
theory in employee development. Includes recommendations for ways in which to
integrate adult learning theory into employee development.

\hypertarget{scholarly-integration-1}{%
\subsubsection*{SCHOLARLY INTEGRATION}\label{scholarly-integration-1}}
\addcontentsline{toc}{subsubsection}{SCHOLARLY INTEGRATION}

\textbf{Unsatisfactory:} Does not integrate scholarly references to support claims and assertions made in
the paper.

\textbf{Developing:} Integrates scholarly references to support some of the claims and assertions
made in the paper.

\textbf{Proficient:} \emph{Integrates scholarly references to support claims and assertions made in the
paper.}

\textbf{Exemplary:} Integrates scholarly references to support claims and assertions made in the
paper, effectively synthesizing different perspectives and research results from
scholarly sources.

\begin{longtable}[]{@{}
  >{\raggedright\arraybackslash}p{(\columnwidth - 8\tabcolsep) * \real{0.2000}}
  >{\raggedright\arraybackslash}p{(\columnwidth - 8\tabcolsep) * \real{0.2000}}
  >{\raggedright\arraybackslash}p{(\columnwidth - 8\tabcolsep) * \real{0.2000}}
  >{\raggedright\arraybackslash}p{(\columnwidth - 8\tabcolsep) * \real{0.2000}}
  >{\raggedright\arraybackslash}p{(\columnwidth - 8\tabcolsep) * \real{0.2000}}@{}}
\toprule\noalign{}
\begin{minipage}[b]{\linewidth}\raggedright
\textbf{TOTAL}
\end{minipage} & \begin{minipage}[b]{\linewidth}\raggedright
\textbf{0 = 0\% (F)}
\end{minipage} & \begin{minipage}[b]{\linewidth}\raggedright
\textbf{10 = 50\% (C)}
\end{minipage} & \begin{minipage}[b]{\linewidth}\raggedright
\textbf{15 = 75 (B)}
\end{minipage} & \begin{minipage}[b]{\linewidth}\raggedright
\textbf{20 = 100\% (A+)}
\end{minipage} \\
\midrule\noalign{}
\endhead
\bottomrule\noalign{}
\endlastfoot
\end{longtable}

\begin{center}\rule{0.5\linewidth}{0.5pt}\end{center}

\hypertarget{assignment-platform-paper}{%
\section*{Assignment: Platform Paper}\label{assignment-platform-paper}}
\addcontentsline{toc}{section}{Assignment: Platform Paper}

\begin{assessment}
For this assignment, you will write a contextualized Platform Paper in
which you discuss your ideal learning community and your role as
teacher/leader of that learning community. Select a context for your
paper (i.e.~facilitating in a FAR Centre in a specific country, teaching
adult learners, facilitating employee development workshops, etc.). Your
paper should be written and referenced in APA format and include
references to a minimum of 10 scholarly sources (this can include
literature you read in this course). You will write a draft of the
Platform Paper in Unit 8 and post for Peer Review. In Unit 9, you will
provide feedback to another learner on their paper. You will make
revisions based on the Peer Review and, in Unit 10, you will submit the
final Platform Paper. Peer reviewers will be assigned in advance.

{Paper Outline}

This paper will be 12-15 pages long, and should include:

\begin{enumerate}
\def\labelenumi{\arabic{enumi}.}
\tightlist
\item
  Introduction (1-2 pages)\\
\item
  Section 1: Ideal Learning Environment (5-7 pages)\\
\item
  Section 2: Your Role as Teacher and Leader (5-7 pages)\\
\item
  Conclusion (1-2 pages)
\end{enumerate}

{Paper Guidelines}

\begin{itemize}
\tightlist
\item
  \textbf{Introduction}: Introduce the two sections in your paper,
  providing a brief description of the key points you will make in each
  section.\\
\item
  \textbf{Section 1}: In section one, you will describe your ideal
  education learning environment. This section should demonstrate your
  learning about authentic learning communities, incorporating scholarly
  sources and your own analysis to depict your ideal learning
  environment. Incorporate a discussion of the learning community
  environment, learning experiences, student learning outcomes, and
  personal beliefs about teaching and learning.\\
\item
  \textbf{Section 2}: In this section, describe your role as a teacher
  or leader within an authentic learning community. Incorporating
  scholarly literature, analyze your role as a facilitator/leader in
  planning learning experiences, facilitating student learning, and
  assessing student learning. Describe the actions, practices, and
  strategies you will engage in to achieve your vision of the learning
  community you described in section one.\\
\item
  \textbf{Conclusion}: Summarize the key points you made in each
  section.\\
\item
  \textbf{References}: Include a reference list with references to at
  least 10 scholarly sources.
\end{itemize}
\end{assessment}

\hypertarget{platform-paper-rubric}{%
\subsection*{Platform Paper Rubric}\label{platform-paper-rubric}}
\addcontentsline{toc}{subsection}{Platform Paper Rubric}

See the following rubric that explains how your assignment will be evaluated. Also available as a \href{assets/assessment/Platform-Paper-RUBRIC.pdf}{pdf}

\hypertarget{apawriting}{%
\subsubsection*{APA/WRITING}\label{apawriting}}
\addcontentsline{toc}{subsubsection}{APA/WRITING}

\textbf{Unsatisfactory:} Paper does not model language and conventions used in scholarly literature. Writing is not well-organized. Several errors in grammar or composition. Sources are not cited. APA citations are not appropriately formatted.

\textbf{Developing:} Paper partially models language and conventions used in scholarly literature. Writing is somewhat well organized and includes some errors in grammar or composition. Not all sources cited. APA citations are generally formatted correctly, with several errors.

\textbf{Proficient:} \emph{Paper consistently models language and conventions used in scholarly literature. Writing is well-organized and includes few (if any) errors in grammar or composition. All resources are appropriately cited (including in-text citations and bibliography information). Few (if any) errors in APA citations.}

\textbf{Exemplary:} Paper is an exemplar of language and conventions used in scholarly literature. Writing is well-organized and free of errors in grammar or composition. All resources are appropriately cited. No errors in APA format.

\hypertarget{developing-a-cohesive-and-logical-academic-argument-2}{%
\subsubsection*{DEVELOPING a COHESIVE and LOGICAL ACADEMIC ARGUMENT}\label{developing-a-cohesive-and-logical-academic-argument-2}}
\addcontentsline{toc}{subsubsection}{DEVELOPING a COHESIVE and LOGICAL ACADEMIC ARGUMENT}

\textbf{Unsatisfactory:} Does not make a focused, cohesive, or logical academic argument. Paper is confusing, and is missing an introduction, body, or conclusion. Transitions between sections and ideas are missing.

\textbf{Developing:} Makes an academic argument that is only partially focused, cohesive and logical. Paper is generally organized, but is missing an introduction, body, or conclusion. Transitions between sections and ideas are unclear.

\textbf{Proficient:} \emph{Makes a focused, cohesive, logical academic argument. Paper is effectively organized and includes an introduction, body, and conclusion. Transitions between sections and ideas are clear.}

\textbf{Exemplary:} Makes a focused, cohesive, logical and compelling academic argument. Paper is effectively organized and includes an introduction, body, and conclusion. Transitions between sections and ideas are clear, and build on each other.

\hypertarget{ideal-learning-environment}{%
\subsubsection*{IDEAL LEARNING ENVIRONMENT}\label{ideal-learning-environment}}
\addcontentsline{toc}{subsubsection}{IDEAL LEARNING ENVIRONMENT}

\textbf{Unsatisfactory:} Does not include a description of your ideal learning environment. Does not reference scholarly sources. Does note analyze key elements of an authentic learning community. Does not mention or describe the learning community environment, student learning outcomes, learning outcomes and personal beliefs about teaching and learning.

\textbf{Developing:} Includes a partial description of your ideal learning environment, referencing few scholarly sources and including a partial analysis of key elements of an authentic learning community. Mentions some elements, but does not fully describe the learning community environment, student learning outcomes, learning outcomes and personal beliefs about teaching and learning.

\textbf{Proficient:} \emph{Includes a detailed description of your ideal learning environment, referencing scholarly sources and analyzing key elements of an authentic learning community. Describes the learning community environment, student learning outcomes, learning outcomes and personal beliefs about teaching and learning.}

\textbf{Exemplary:} Includes a detailed description of your ideal learning environment, referencing scholarly sources and analyzing key elements of authentic learning communities. Provides a rationale for key elements of the learning community environment, student learning outcomes, learning outcomes and personal beliefs about teaching and learning. Advances scholarly thinking about authentic learning communities.

\hypertarget{your-role-as-teacher-and-leaders}{%
\subsubsection*{YOUR ROLE AS TEACHER AND LEADERS}\label{your-role-as-teacher-and-leaders}}
\addcontentsline{toc}{subsubsection}{YOUR ROLE AS TEACHER AND LEADERS}

\textbf{Unsatisfactory:} Does not include a description of your role as a teacher or leader within an authentic learning community, incorporating scholarly literature. Does not include an analysis of your role as a facilitator/leader in planning learning experiences, facilitating student learning, and assessing student learning. Does not include a description of the actions, practices, and strategies you will engage in to achieve your vision of the learning community you described in section one.

\textbf{Developing:} Includes a partial description of your role as a teacher or leader within an authentic learning community, incorporating scholarly literature. Describes but does not analyze your role as a facilitator/leader in planning learning experiences, facilitating student learning, and assessing student learning. Lists but does not describe the actions, practices, and strategies you will engage in to achieve your vision of the learning community you described in section one.

\textbf{Proficient:} \emph{Includes a detailed description of your role as a teacher or leader within an authentic learning community, incorporating scholarly literature. Includes a detailed analysis of your role as a facilitator/leader in planning learning experiences, facilitating student learning, and assessing student learning. Includes a detailed description of the actions, practices, and strategies you will engage in to achieve your vision of the learning community you described in section one.}

\textbf{Exemplary:} Includes a detailed analysis of your role as a teacher or leader within an authentic learning community, incorporating scholarly literature. Includes a detailed analysis of your role as a facilitator/leader in planning learning experiences, facilitating student learning, and assessing student learning. Includes a detailed description of the actions, practices, and strategies you will engage in to achieve your vision of the learning community you described in section one. Synthesizes scholarly thinking about the role of the teacher/leader.

\hypertarget{scholarly-integration-2}{%
\subsubsection*{SCHOLARLY INTEGRATION}\label{scholarly-integration-2}}
\addcontentsline{toc}{subsubsection}{SCHOLARLY INTEGRATION}

\textbf{Unsatisfactory:} Does not integrate many references to support the arguments made in the paper.

\textbf{Developing:} Integrates fewer than 10 scholarly sources to support arguments made in the paper.

\textbf{Proficient:} \emph{Integrates a minimum of 10 scholarly sources to support arguments made in each section of the paper.}

\textbf{Exemplary:} Integrates a minimum of 10 references to support the arguments made in each section, including several scholarly sources not included in course materials.

\begin{longtable}[]{@{}
  >{\raggedright\arraybackslash}p{(\columnwidth - 8\tabcolsep) * \real{0.2000}}
  >{\raggedright\arraybackslash}p{(\columnwidth - 8\tabcolsep) * \real{0.2000}}
  >{\raggedright\arraybackslash}p{(\columnwidth - 8\tabcolsep) * \real{0.2000}}
  >{\raggedright\arraybackslash}p{(\columnwidth - 8\tabcolsep) * \real{0.2000}}
  >{\raggedright\arraybackslash}p{(\columnwidth - 8\tabcolsep) * \real{0.2000}}@{}}
\toprule\noalign{}
\begin{minipage}[b]{\linewidth}\raggedright
\textbf{TOTAL}
\end{minipage} & \begin{minipage}[b]{\linewidth}\raggedright
\textbf{0 = 0\% (F)}
\end{minipage} & \begin{minipage}[b]{\linewidth}\raggedright
\textbf{10 = 50\% (C)}
\end{minipage} & \begin{minipage}[b]{\linewidth}\raggedright
\textbf{15 = 75 (B)}
\end{minipage} & \begin{minipage}[b]{\linewidth}\raggedright
\textbf{20 = 100\% (A+)}
\end{minipage} \\
\midrule\noalign{}
\endhead
\bottomrule\noalign{}
\endlastfoot
\end{longtable}

\hypertarget{references}{%
\chapter*{References}\label{references}}
\addcontentsline{toc}{chapter}{References}

The following are key references used in this course. \textbf{\emph{Check with your course syllabus for required readings.}}

  \bibliography{book.bib}

\end{document}
